\chapter{Discussion}\label{ch:discussion}

\section{Interpretation}
The attribution of \S\ref{sec:attribution} assigns the two halves of the
result to two parts of the model: per-piece feature weighting recovers the
hidden values, and the score graph calibrates the uncertainty. The posterior
decomposition sharpens this picture from inside a single fit. The graph
component and the music-model component act nearly orthogonally at each
note, with a component correlation of $\approx 0$, so the two are complements
rather than rivals. Which of them dominates is itself a per-piece property
that the evidence assigns, and the assignment can invert completely between
pieces (Fig.~\ref{fig:component-contrast}). Two questions remain that the
numbers alone do not answer: why structure improves calibration rather than
only accuracy, and why the music model's marginal contribution is
concentrated in a single channel.

Structure improves calibration because coupling correlated notes lets each
held-out note borrow statistical strength from its neighbours. Borrowing
sharpens the posterior mean, but its more important effect is on the
posterior variance, which then reflects how well a note is actually
determined by its context. An independent predictor must choose a single
global uncertainty for every note. The graph gives each note its own, and
that per-note uncertainty is what brings coverage to the nominal level.

The music model's marginal contribution is concentrated in velocity because
dynamics is the channel most predictable from learned performance
conventions. Phrase-level swells and metrical stress are conventions that a
sequence model captures and that local hand-built features do not. Timing is
the channel on which every model's error plateaus, since much of its
variance is note-local. Articulation is comparatively idiosyncratic. On
those two channels the music model's additional information therefore has
the least to add.

\section{Limitations}
Two results are supported on held-out test sets that were evaluated once
against hypotheses fixed in advance (\S\ref{sec:hygiene}): the symbolic
piano result and the six-channel result on real recordings. The latter rests
on a single corpus, the 13 held-out pieces of URMP, and two of its findings
run against the model. The secondary timing-calibration hypothesis was not
supported, and timing recovery showed a small adverse difference whose
interval excludes zero. Two follow-up hypotheses
were then evaluated on the same strings-and-winds test set a second time,
with the re-use disclosed (\S\ref{sec:phase2-reuse}). The revised timing
noise model was supported, and the tonal metric's primary calibration
prediction was not. Evidence from a re-used test set is weaker in kind than
evidence from a test set evaluated once.

The waveform evidence is exploratory and comes entirely from the development
set. Its tests of the thesis's methodological hypothesis, three spans
covering short chains, note joints and phrases
(\S\ref{sec:phase3-priorvalue}), support only the weak form of that
hypothesis. The score prior outperforms a purely temporal prior at phrase
scale, but beyond the context's common offset it adds nothing measurable at
the spans tested. A further development-set finding qualifies the graph
itself. On both corpora the graph's value is reproduced, and on piano
slightly exceeded, by the same graph with its pitch term removed
(Appendix~\ref{app:trackscale}). The structure that does the work is
therefore smoothness in score time over about two beats, not the pitch
coupling that makes the graph a score graph. The results supported on the
held-out test sets stand as comparisons of the full model with its
alternatives, but their attribution is narrower than the model's name
suggests.

The remaining limitations concern the data and the model's form. The piano
targets depend on the quality of a single corpus's
score$\leftrightarrow$performance alignment (ASAP), and the evaluation sets,
although held out and filtered for contamination, are modest in size, with
30 development pieces and 20 held-out test pieces. Two model-level
limitations are observed rather than hypothetical. The Gaussian likelihood's
quadratic tail can amplify a handful of timing outliers on a very steady
piece into the pooled NLL (\S\ref{sec:guard-gp}), and per-piece feature
weighting breaks down under excerpt extrapolation (\S\ref{sec:boundary}).
Generalisation beyond these settings is untested.

\section{Threats to validity}
The main threats to validity were addressed explicitly rather than assumed
away. Target leakage was found and corrected, and the correction was
verified by structural tests and by a bitwise end-to-end audit of the
proposed model's path (\S\ref{ch:results}). Train/test contamination was
filtered (\S\ref{sec:contam}). Hyperparameter fit collapse was safeguarded
against and made visible by reporting the median and the worst piece
(\S\ref{sec:eb}). Multiple-comparison and selection pressure was met by
per-piece paired significance throughout and, decisively, by the held-out
test set of \S\ref{sec:hygiene}, evaluated once against hypotheses fixed in
advance, which exposed one development-set advantage that did not replicate.
The threat that remains open is external validity beyond the two corpora
studied, ASAP for piano and URMP for chamber ensembles.

% =====================================================================
